%! Properties and Applications of Determinants — Unit 5, Lesson 5.2 — Lesson Plan
\documentclass[10pt]{article}
\usepackage{linalg-article}
\usepackage{linalg-boxes}
\usepackage{graphicx}   % -article does NOT load graphicx; needed for thumbnails/screenshots
\graphicspath{{images/}}
\newcommand{\TallMath}[1]{$\displaystyle #1\rule[-1.4em]{0pt}{3.2em}$}
\newcommand{\avec}{\mathbf{a}}
\newcommand{\bb}{\mathbf{b}}
\newcommand{\xx}{\mathbf{x}}
\newcommand{\T}{\mathsf{T}}
\newcommand{\UnitNumberName}{Unit 5: Determinants and Linear Transformations \quad}
\newcommand{\LessonNumberName}{Lesson 5.2: Properties and Applications of Determinants}

\begin{document}
\begin{center}
    {\Large\bfseries \CourseName: \SchoolYear \\
    {\small\bfseries \UnitNumberName \LessonNumberName}}
\end{center}

\begin{tcolorbox}[colback=blush, colframe=burgundy, boxrule=0.9pt, arc=2mm,
    left=2mm, right=2mm, top=1.5mm, bottom=1.5mm]
    \textbf{Primary Objective:} Students can use the \emph{properties} of determinants --- read
    \emph{geometrically} as facts about volume --- instead of always expanding: (1) the three founding
    rules ($\det I=1$; a \emph{swap} of two rows/columns \emph{negates} the determinant; scaling one
    row/column by $t$ scales $\det$ by $t$); (2) the key property that \emph{adding a multiple of one
    row to another} (a \emph{shear}) leaves $\det$ \emph{unchanged} --- so elimination barely touches
    it. They can therefore compute a determinant the \emph{fast} way --- reduce to triangular $U$ and
    take $\det A=\pm(\text{product of the pivots})$ --- and use the \emph{product rule}
    $\det(AB)=\det A\,\det B$ and its consequences ($\det A^{-1}=1/\det A$, $\det A^{\T}=\det A$,
    and $\det A=0\Leftrightarrow$ singular). They can say where Unit~5 goes next: \emph{linear
    transformations} (5.3), where $|\det A|$ is the volume-scaling factor.
\end{tcolorbox}

\renewcommand{\arraystretch}{1.4}
\begin{skillbox}[Priority Ideas \& Skills]{goldbox}
\begin{tabularx}{\linewidth}{@{}X|X@{}}
\textbf{Priority Ideas \& Skills} & \textbf{Key Understandings (the ``why'')} \\[2pt]
\hline
\vspace{-6pt}
\begin{itemize}[leftmargin=1.1em, nosep, after=\vspace{2pt}]
  \item Know the founding rules: $\det I=1$; \textbf{swap} negates; scaling a row by $t$ scales $\det$ by $t$.
  \item Use the key property: \textbf{adding a multiple of one row to another leaves $\det$ unchanged}.
  \item Compute $\det$ the fast way: reduce to triangular, $\det=\pm(\text{product of the pivots})$.
  \item Use $\det(AB)=\det A\,\det B$, so $\det A^{-1}=1/\det A$ and $\det A=0\Leftrightarrow$ singular.
\end{itemize}
&
\vspace{-6pt}
\begin{itemize}[leftmargin=1.1em, nosep, after=\vspace{2pt}]
  \item Every rule is a fact about \emph{volume}: swap $=$ reflect, scale $=$ stretch, shear $=$ slide.
  \item A \emph{shear} keeps base and height, so it keeps volume --- that is why elimination preserves $\det$.
  \item Elimination already makes matrices triangular; the pivots hand you the determinant for free.
  \item Composing maps multiplies volume factors --- so determinants multiply too.
\end{itemize}
\end{tabularx}
\end{skillbox}

\begin{skillbox}[{Vocabulary, Concepts, \& Theorems}]{greenbox}
\renewcommand{\arraystretch}{1.3}
\begin{tabularx}{\linewidth}{@{}l X@{}}
\textbf{Founding rules} & $\det I=1$; swapping two rows (or columns) reverses the sign; scaling a row by $t$ scales $\det$ by $t$. \\
\textbf{Shear invariance} & Adding a multiple of one row to another does \emph{not} change $\det$ (the elimination step). \\
\textbf{Product of pivots} & Reduce $A$ to triangular $U$ by elimination; $\det A=\pm(\text{product of the pivots})$, one sign flip per row swap. \\
\textbf{Product rule} & $\det(AB)=\det A\,\det B$; hence $\det A^{-1}=1/\det A$ and $\det A^{\T}=\det A$. \\
\textbf{Cramer's rule} (application) & Solve $A\xx=\bb$ with ratios of determinants: $x_i=\det(A_i)/\det A$. \\
\end{tabularx}
\end{skillbox}

\begin{fixedskillbox}[Activate Prior Knowledge \& Spiral Review]{blush}
    The Warm-Up rehearses the three moves this lesson fuses into one shortcut, all on the same tiny
    matrix $M=\begin{bsmallmatrix} 1 & 2\\ 2 & 5 \end{bsmallmatrix}$:
    \textbf{(1)} compute a \emph{$2\times2$ determinant} ($\det M=1$, Lesson~5.0);
    \textbf{(2)} take \emph{one elimination step} (Unit~2) to reach triangular $U=\begin{bsmallmatrix} 1 & 2\\ 0 & 1 \end{bsmallmatrix}$;
    and \textbf{(3)} use the \emph{triangular shortcut} (Lesson~5.1) to get $\det U=1\cdot1=1$. Debrief the
    punchline aloud: $\det M=\det U$ --- \emph{elimination did not change the determinant.} That single
    observation is the engine of the whole lesson.
\end{fixedskillbox}

\begin{skillbox}[Hook]{blush}
    Open on the pain point students just felt: expanding a $3\times3$ by cofactors is \emph{three}
    $2\times2$ determinants and a lot of bookkeeping. Ask: \textbf{``Is there a faster way?''} Point out
    that they \emph{already} own a machine that turns any matrix triangular --- \textbf{elimination}
    from Unit~2 --- and a triangular determinant is free (just multiply the diagonal, Lesson~5.1). The
    only missing link: \textbf{does elimination change the determinant?} Send them to the Warm-Up to
    find out on a $2\times2$; the answer (\emph{no}) is the key that unlocks ``$\det=$ product of the
    pivots.'' Geometric punchline to hold up: an elimination step is a \textbf{shear} --- it slides the
    box sideways but keeps its base and height, so the \emph{volume never changes}.
\end{skillbox}

\begin{skillbox}[Lesson]{blush}
\begin{multicols}{2}
\textbf{1. The rules determinants live by.} Three founding facts, each a statement about the box's
volume: $\det I=1$ (the unit cube has volume $1$); \emph{swapping} two rows/columns \textbf{negates}
$\det$ (a mirror reflection --- same size, flipped orientation); scaling one row by $t$ scales $\det$
by $t$ (stretch one edge, stretch the volume). A quick consequence: two \emph{equal} rows $\Rightarrow$
$\det=0$ (a squashed, flat box).

\textbf{2. The key property: shear invariance.} \emph{Adding a multiple of one row to another does not
change $\det$.} Geometrically it is a \textbf{shear} --- slide the top of the box sideways; base and
height are untouched, so volume is unchanged. This is \emph{exactly} the elimination step from Unit~2.
Show it on $\begin{bsmallmatrix} 2 & 1\\ 1 & 3 \end{bsmallmatrix}$ (still $\det=5$ after a shear).

\columnbreak

\textbf{3. The fast way: $\det=\pm$(product of pivots).} Since shears don't change $\det$ and swaps
only flip its sign, run elimination to triangular $U$ and \emph{multiply the pivots}. Worked:
$\begin{bsmallmatrix} 1 & 2 & 0\\ 2 & 5 & 1\\ 0 & 1 & 3 \end{bsmallmatrix}\to$ pivots $1,1,2\to\det=2$
--- no cofactors. (Count row swaps: each one flips the sign.)

\textbf{4. The product rule \& applications.} $\det(AB)=\det A\,\det B$ (compose two maps $\Rightarrow$
multiply the volume factors). Worked: $\det A=5$, $\det B=6\Rightarrow\det(AB)=30$. Consequences:
$\det A^{-1}=1/\det A$, $\det A^{\T}=\det A$, and $\det A=0\Leftrightarrow$ \emph{singular}. \textbf{Road
ahead:} \textbf{5.3} linear transformations, where $|\det A|$ is the volume-scaling factor.
\end{multicols}
\end{skillbox}

\begin{skillbox}[Active Monitoring]{redbox}
Circulate during the notes practice and Tier work. Watch for and correct:
\begin{itemize}[leftmargin=1.2em, nosep]
  \item confusing the row rules --- a \emph{swap} negates and a \emph{scale} multiplies, but ``add a multiple of one row to another'' changes \emph{nothing};
  \item forgetting the \emph{sign flip} per row swap when reading $\det$ off the pivots;
  \item thinking $\det(A+B)=\det A+\det B$ (false) --- the product rule is about $AB$, not $A+B$.
\end{itemize}
Cold-call prompts: ``You added $2\times$row~1 to row~2 --- what happened to the determinant?'' ``How
many swaps did you make, and what does that do to the sign?'' ``$\det A=5$ --- what is $\det A^{-1}$?''
\end{skillbox}

\begin{skillbox}[Group Work \& Differentiation]{redbox}
\textbf{Rules, Pivots, and Products} activity (tiered; mirrors the notes).
\begin{multicols}{3}
\textbf{Tier R --- Remediate.} Compute a determinant by reducing to triangular (product of pivots);
check that a shear leaves $\det$ unchanged; see that a repeated column forces $\det=0$.

\columnbreak
\textbf{Tier A --- Approaching.} Reduce with a \emph{swap} and track the sign; verify the product rule
$\det(AB)=\det A\det B$; use $\det A^{-1}=1/\det A$ and $\det A=0\Rightarrow$ no inverse.

\columnbreak
\textbf{Tier E --- Extension.} \emph{An application:} \textbf{Cramer's rule} --- solve a $2\times2$
system with ratios of determinants --- and confirm $\det A^{\T}=\det A$.
\end{multicols}
\end{skillbox}

\begin{skillbox}[Individual Work \& Assessment]{redbox}
\textbf{Exit Ticket} (independent, no notes): compute a $3\times3$ determinant \emph{by elimination}
(product of the pivots); use the \emph{product rule} to get $\det(AB)$ and $\det A^{-1}$ from
$\det A,\det B$; and a \textbf{justification check} --- why does adding a multiple of one row to
another leave the determinant unchanged? Collect and sort into ``got it / rule slip / sign slip'' to
decide whether 5.3 opens with a quick properties re-warm.
\end{skillbox}

\begin{skillbox}[Reinforcement \& Extension]{goldbox}
\textbf{Homework:} compute determinants by elimination (including one triangular and one with a
dependent row $\Rightarrow0$); apply the row rules (swap, scale, shear); use the product rule and
$\det A^{-1}=1/\det A$; and a short justification item. \textbf{Extension:} \emph{Cramer's rule} for a
$2\times2$ system, plus a check that $\det A^{\T}=\det A$. \textbf{Preview:} Lesson~5.3, \emph{Linear
Transformations} --- where $|\det A|$ is the factor every area/volume is multiplied by, and the sign
tells you whether orientation was preserved or flipped.
\end{skillbox}
\end{document}
